\documentclass[11pt,a4paper]{article}

% --- Required Packages ---
\usepackage[margin=1in]{geometry}
\usepackage{amsmath,amssymb,amsfonts,bm,mathtools,amsthm}
\usepackage{booktabs}
\usepackage[hidelinks]{hyperref} % Added [hidelinks] to remove red boxes
\usepackage{bookmark} % Added to prevent rerunfilecheck outline warnings

% --- Theorem Environments ---
\newtheorem{theorem}{Theorem}

% --- Document Metadata ---
\title{\textbf{Diagnostic Classification Models: Rigorous Mathematical 
Formulation of DINA, DINO, rDINA, and rDINO Models 
with Prior Specifications}}
\author{\textbf{Jonathon A. Pedroza \& Boting Zhang}}
\date{\today}

\begin{document}

\maketitle

% --- Table of Contents ---
\tableofcontents
\newpage

\section{Latent Space Formulation}

Let $N$ denote the total number of examinees ($i = 1, \dots, N$), 
$J$ denote the total number of items ($j = 1, \dots, J$), 
and $K$ denote the number of binary latent attributes ($k = 1, \dots, K$). 
In the provided Stan models, $J$ represents examinees, 
$I$ represents items, and $C$ represents the latent classes. 
For notational standard, we map these to $N$ (examinees), 
$J$ (items), and $L$ (latent classes).

\begin{itemize}
    \item \textbf{Latent Attribute Vector:} 
    Examinee $i$'s attribute mastery profile is represented by the binary vector:
    \begin{equation}
        \boldsymbol{\alpha}_i = (\alpha_{i1}, \alpha_{i2}, \dots, \alpha_{iK})^T \in \{0, 1\}^K
    \end{equation}
    where $\alpha_{ik} = 1$ if examinee $i$ has mastered attribute $k$, 
    and $0$ otherwise.

    \item \textbf{Latent Class Space:} 
    The latent class space comprises $L = 2^K$ distinct attribute patterns. 
    Let $\boldsymbol{\alpha}_l$ denote the $l$-th attribute profile vector 
    for $l \in \{1, \dots, L\}$. 
    The prior distribution of latent class membership is governed 
    by the structural probability vector:
    \begin{equation}
        \boldsymbol{\pi} = (\pi_1, \pi_2, \dots, \pi_L)^T, 
        \quad \text{where } \pi_l = P(\boldsymbol{\alpha}_i = \boldsymbol{\alpha}_l), 
        \quad \sum_{l=1}^L \pi_l = 1
    \end{equation}

    \item \textbf{Q-Matrix ($\mathbf{Q}$):} 
    $\mathbf{Q} = [q_{jk}]_{J \times K} \in \{0, 1\}^{J \times K}$, 
    where $q_{jk} = 1$ if item $j$ requires attribute $k$, 
    and $0$ otherwise.
\end{itemize}

% Fixed hyperref warning using \texorpdfstring
\section{Deterministic Inputs and Ideal Responses (\texorpdfstring{$\xi_{jl}$}{xi\_jl})}

The core difference between the conjunctive (DINA) and disjunctive (DINO) 
families lies in the formulation of the deterministic input (ideal response), 
denoted as $\xi_{jl}$ (or $\eta_{jl}$), which represents whether 
latent class $l$ meets the attribute requirements for item $j$.

\subsection{DINA (Deterministic Inputs, Noisy ``AND'' Gate)}
In the DINA model, an examinee must possess \textit{all} required attributes 
for an item to have an ideal response of 1. 
It functions as a logical AND gate:
\begin{equation}
    \xi_{jl}^{\text{DINA}} = \prod_{k=1}^K \alpha_{lk}^{q_{jk}} \in \{0, 1\}
\end{equation}

\subsection{DINO (Deterministic Inputs, Noisy ``OR'' Gate)}
In the DINO model, an examinee must possess \textit{at least one} 
of the required attributes to have an ideal response of 1. 
It functions as a logical OR gate:
\begin{equation}
    \xi_{jl}^{\text{DINO}} = 1 - \prod_{k=1}^K (1 - \alpha_{lk})^{q_{jk}} \in \{0, 1\}
\end{equation}

\section{Item Response Functions and Stochastic Noise Models}

To account for stochastic noise (i.e., human error), 
we apply probabilistic parameters to the deterministic inputs.

\subsection{Unrestricted Models: DINA and DINO}
The standard DINA and DINO models utilize two item-specific noise parameters:
\begin{enumerate}
    \item \textbf{Slipping Parameter ($s_j$):} $s_j = P(Y_{ij} = 0 \mid \xi_{jl} = 1)$. 
    The probability of missing the item despite meeting attribute requirements.
    \item \textbf{Guessing Parameter ($g_j$):} $g_j = P(Y_{ij} = 1 \mid \xi_{jl} = 0)$. 
    The probability of correctly guessing the item without meeting requirements.
\end{enumerate}
The conditional success probability for latent class $l$ on item $j$ is formulated as:
\begin{equation}
    P_{jl} \equiv P(Y_{ij} = 1 \mid \boldsymbol{\alpha}_l, s_j, g_j) 
    = (1 - s_j)^{\xi_{jl}} g_j^{1 - \xi_{jl}}
\end{equation}

\subsection{Restricted Models: rDINA and rDINO}
The restricted models (rDINA and rDINO) impose a parsimony constraint 
by assuming that the probability of slipping is equal to 
the probability of guessing. They define a single noise parameter per item:
\begin{equation}
    e_j = s_j = g_j
\end{equation}
The conditional success probability is simplified to:
\begin{equation}
    P_{jl} \equiv P(Y_{ij} = 1 \mid \boldsymbol{\alpha}_l, e_j) 
    = (1 - e_j)^{\xi_{jl}} e_j^{1 - \xi_{jl}}
\end{equation}
This formulation restricts the flexibility of the item response function 
but enhances parameter identifiability and reduces standard errors, 
especially in small sample sizes.

\section{Likelihood Structure and Prior Formulations}

\subsection{Marginal Log-Likelihood}
For an observed response matrix $\mathbf{Y}$ of $N$ respondents and $J$ items, 
the marginal log-likelihood, integrating over the $L$ discrete latent classes, 
is identical in structure across the DINA, DINO, rDINA, and rDINO models. 
The formulation models the probability of the observed response $Y_{ij} \in \{0, 1\}$ 
given the item parameters (slipping $\mathbf{s}$, guessing $\mathbf{g}$) 
and latent class proportions $\boldsymbol{\pi}$:

$$ \log L(\mathbf{s}, \mathbf{g}, \boldsymbol{\pi} \mid \mathbf{Y}) 
= \sum_{i=1}^N \log \left( \sum_{l=1}^L \pi_l 
\prod_{j=1}^J P_{jl}^{Y_{ij}} (1 - P_{jl})^{1 - Y_{ij}} \right) $$

where $P_{jl}$ represents the probability of a respondent in latent class $l$ 
correctly answering item $j$. The definition of $P_{jl}$ depends on the 
specific deterministic input (e.g., the $\xi_{jl}$ or $\eta_{jl}$ matrices 
derived from the Q-matrix) dictating whether the class possesses 
the required attributes.

\subsection{Estimation Paradigms and Prior Specifications}
The models demonstrate estimation under three distinct paradigms using Stan: 
Maximum Likelihood Estimation (MLE), Maximum A Posteriori (MAP), 
and fully Bayesian Markov Chain Monte Carlo (MCMC). 
These are controlled by modifying the objective function via 
the `target +=` statement and the chosen estimation algorithm.

\subsubsection{Maximum Likelihood Estimation (MLE / Flat Posterior)}
When optimized using the L-BFGS algorithm without explicit prior distributions, 
the model implies $s_j, g_j \sim \text{Beta}(1, 1)$. 
This prior acts as a uniform distribution over $[0, 1]$. 
Because the log-density of a $\text{Beta}(1,1)$ distribution is $0$, 
it adds a constant to the log-posterior. The target corresponds purely to the MLE:

$$ \log P(\boldsymbol{\theta} \mid \mathbf{Y}) 
\propto \log L(\mathbf{s}, \mathbf{g}, \boldsymbol{\pi} \mid \mathbf{Y}) $$

While computationally efficient, MLE point estimates in sparse diagnostic data 
often drift toward boundaries ($0$ or $1$). 
Mathematically, this forces reliability metrics like accuracy ($\max(p, 1-p)$) 
and consistency ($p^2 + (1-p)^2$) to artificially resolve to exactly $1.000$.

\subsubsection{Maximum A Posteriori (MAP / Regularization)}
To regularize the model and prevent boundary estimates (e.g., $s_j=0$ or $g_j=0$), 
mildly informative priors are applied and optimized via L-BFGS. 
Specifically, $s_j, g_j \sim \text{Beta}(5, 20)$. 
This reflects prior substantive knowledge that slipping and guessing 
parameters generally fall below $0.30$, with an expected mean of $5 / (5+20) = 0.20$.

The regularized log-posterior (penalized log-likelihood) becomes:

$$ \log P(\boldsymbol{\theta} \mid \mathbf{Y}) 
\propto \log L(\boldsymbol{\theta} \mid \mathbf{Y}) 
+ \sum_{j=1}^J \Big[ (5 - 1)\log(s_j) + (20 - 1)\log(1 - s_j) \Big] + \dots $$

\begin{theorem}[Well-Posedness of Regularized Optimization]\label{thm:well_posed}
Under the $\text{Beta}(5, 20)$ prior, the shape parameters strictly exceed $1$. 
This ensures that the penalization terms $\log(s_j)$ and $\log(1-s_j)$ 
force the objective function toward $-\infty$ as $s_j \to 0$ or $s_j \to 1$. 
Thus, the MAP estimates are strictly constrained to the interior of the 
parameter space $(0, 1)$, resolving numerical instabilities associated 
with zero-counts.
\end{theorem}
\textit{Proof provided in Appendix \ref{app:proof_well_posed}.}

\subsubsection{Markov Chain Monte Carlo (MCMC) Estimation}
For full uncertainty quantification, the models are additionally estimated via MCMC 
utilizing Stan's Hamiltonian Monte Carlo (HMC) with the No-U-Turn Sampler (NUTS). 
Unlike MLE/MAP which yield point estimates, MCMC yields full posterior 
distributions for $\mathbf{s}$, $\mathbf{g}$, and class memberships. 

MCMC estimation is deployed under two prior conditions:
\begin{itemize}
    \item \textbf{Non-Informative Priors:} Utilizing $\text{Beta}(1, 1)$ priors, 
    allowing the posterior to be entirely data-driven.
    \item \textbf{Slightly Informative Priors:} Utilizing $\text{Beta}(5, 20)$ 
    for item parameters to gently guide the sampler and improve convergence 
    (measured via $\hat{R} < 1.01$ and $\text{ESS} > 400$).
\end{itemize}

The fully Bayesian approach enables robust Posterior Predictive Model Checks (PPMC). 
By generating replicated datasets ($\mathbf{Y}^{\text{rep}}$) from the posterior draws, 
we compute Posterior Predictive P-values (PPP) at both the global and item levels. 
PPP values near $0.5$ indicate the generated data closely matches 
the observed data structure, confirming that the specified Q-matrix and 
DINA/DINO condensation rules adequately capture the students' actual 
response generation process.

\section{Theoretical Properties and Model Fit Diagnostics}

Building upon the estimation paradigms established in the previous section, 
we now formalize the mathematical advantages of our current testing framework. 
Specifically, we demonstrate the regularizing effect of the MAP estimator 
on parameter identifiability and establish the theoretical basis for our 
Posterior Predictive Model Checks (PPMC).

\subsection{Identifiability and Uniqueness in Restricted Models}

As observed in our empirical tests, Maximum Likelihood Estimation (MLE) 
for diagnostic models often suffers from flat likelihood surfaces or 
boundary estimates when sample sizes are small or class distributions 
are highly skewed. The restricted models (rDINA and rDINO) combined with 
Maximum A Posteriori (MAP) estimation provide a mathematically guaranteed 
unique solution under specific conditions.

Let us analyze the complete-data log-likelihood for a single item $j$ 
under the rDINA/rDINO framework, assuming latent class memberships $\boldsymbol{\alpha}_i$ 
are temporarily known (e.g., within an EM algorithm's M-step or a 
Hamiltonian Monte Carlo iteration). 

Let $N_{j1}$ be the number of examinees for whom $\xi_{jl} = 1$, 
and let $S_{j1}$ be the number of correct responses among them. 
Similarly, let $N_{j0}$ be the number of examinees for whom $\xi_{jl} = 0$, 
with $S_{j0}$ correct responses. 
The item-specific complete-data penalized log-likelihood, 
utilizing an $e_j \sim \text{Beta}(a, b)$ prior, is:

\begin{equation}
\begin{split}
    \ell(e_j) &= S_{j1} \log(1 - e_j) + (N_{j1} - S_{j1}) \log(e_j) 
    + S_{j0} \log(e_j) \\
    &\quad + (N_{j0} - S_{j0}) \log(1 - e_j) 
    + (a - 1)\log(e_j) + (b - 1)\log(1 - e_j)
\end{split}
\end{equation}

Grouping the terms by $\log(e_j)$ and $\log(1-e_j)$, 
we define constants $A_j$ and $B_j$:
\begin{align*}
    A_j &= (N_{j1} - S_{j1}) + S_{j0} + (a - 1) \\
    B_j &= S_{j1} + (N_{j0} - S_{j0}) + (b - 1)
\end{align*}
yielding the simplified objective function:
\begin{equation}
    \ell(e_j) = A_j \log(e_j) + B_j \log(1 - e_j)
\end{equation}

\begin{theorem}[Strict Concavity of the rDINA/rDINO Objective Function]\label{thm:strict_concavity}
If the prior shape parameters satisfy $a > 1$ and $b > 1$, 
the penalized complete-data log-likelihood for the noise parameter $e_j$ 
is strictly concave on the interval $(0, 1)$, 
guaranteeing a unique global maximum.
\end{theorem}
\textit{Proof provided in Appendix \ref{app:proof_strict_concavity}.}

This highlights the core improvement in our current MAP/MCMC testing phase: 
by utilizing $a=5$ and $b=20$, we forcefully bound $A_j$ and $B_j$ away from zero, 
entirely eliminating the multimodal or unbounded likelihood spaces 
that typically cause L-BFGS to fail in sparse Q-matrix designs.

\subsection{Improvements in Model Fit Testing via PPMC}

Our current testing phase relies heavily on Posterior Predictive Model 
Checking (PPMC) to evaluate the absolute fit of the Q-matrix and 
the chosen condensation rule (DINA vs. DINO). 
Unlike traditional frequentist $\chi^2$ tests which rely on asymptotic 
approximations that fail under sparse cognitive diagnostic data, 
PPMC provides exact finite-sample validity.

Let $T(\mathbf{Y})$ be a discrepancy measure 
(e.g., item-pair odds ratios or raw score distributions). 
At each MCMC iteration $m = 1, \dots, M$, we draw a replicated dataset 
$\mathbf{Y}^{\text{rep},(m)}$ from the posterior predictive distribution:
\begin{equation}
    P(\mathbf{Y}^{\text{rep}} \mid \mathbf{Y}) 
    = \int P(\mathbf{Y}^{\text{rep}} \mid \boldsymbol{\theta}) 
    P(\boldsymbol{\theta} \mid \mathbf{Y}) d\boldsymbol{\theta}
\end{equation}

The Bayesian Posterior Predictive P-value (PPP) is defined mathematically 
as the expected indicator function over the posterior:
\begin{equation}
\begin{split}
    PPP &= \mathbb{E}_{\boldsymbol{\theta} \mid \mathbf{Y}} 
    \left[ \Pr(T(\mathbf{Y}^{\text{rep}}, \boldsymbol{\theta}) 
    \ge T(\mathbf{Y}, \boldsymbol{\theta}) \mid \boldsymbol{\theta}) \right] \\
    &\approx \frac{1}{M} \sum_{m=1}^M 
    I\left(T(\mathbf{Y}^{\text{rep},(m)}, \boldsymbol{\theta}^{(m)}) 
    \ge T(\mathbf{Y}, \boldsymbol{\theta}^{(m)})\right)
\end{split}
\end{equation}

By transitioning our tests from purely relative information criteria (AIC/DIC) 
to item-level PPP evaluations, we can independently verify the 
$q$-vector for each item. 
If an item exhibits an extreme PPP value (e.g., $p < 0.05$ or $p > 0.95$) 
under the DINA model but resolves near $0.50$ under the DINO model, 
we possess direct mathematical justification that the latent traits 
interact compensatorily rather than conjunctively for that specific item.

\newpage
\appendix
\section{Proofs of Theorems}

\subsection{Proof of Theorem \ref{thm:well_posed}} \label{app:proof_well_posed}
\begin{proof}
Let the objective function for the slipping parameter $s_j$ 
(and symmetrically for $g_j$) be denoted as $f(s_j)$. 
Isolating the terms involving $s_j$, the regularized log-posterior is:
\begin{equation}
    f(s_j) = \ell(s_j \mid \mathbf{Y}, \mathbf{g}, \boldsymbol{\pi}) 
    + (\alpha - 1)\log(s_j) + (\beta - 1)\log(1 - s_j)
\end{equation}
where $\ell(s_j \mid \dots)$ is the marginal log-likelihood component 
associated with $s_j$, and the Beta prior shape parameters are 
$\alpha = 5$ and $\beta = 20$.

By definition, the likelihood function returns a probability bounded by $[0, 1]$. 
Therefore, its logarithm is bounded above such that 
$\ell(s_j \mid \dots) \le 0$ for all $s_j \in [0, 1]$.

We analyze the limits of the objective function $f(s_j)$ 
as the parameter approaches the boundaries of its domain:
\begin{itemize}
    \item As $s_j \to 0^+$, the penalization term 
    $(\alpha - 1)\log(s_j) = 4\log(s_j) \to -\infty$.
    \item As $s_j \to 1^-$, the penalization term 
    $(\beta - 1)\log(1 - s_j) = 19\log(1 - s_j) \to -\infty$.
\end{itemize}

Because the log-likelihood term $\ell(s_j \mid \dots)$ is bounded 
above by $0$ and cannot diverge to $+\infty$ to cancel out the penalty, 
the entire objective function diverges to negative infinity at the boundaries:
\begin{equation}
    \lim_{s_j \to 0^+} f(s_j) = -\infty \quad \text{and} \quad \lim_{s_j \to 1^-} f(s_j) = -\infty
\end{equation}

Since $f(s_j)$ is a continuous function on the open interval $(0, 1)$ 
and tends to $-\infty$ at both boundaries, by the Extreme Value Theorem 
(adapted for open intervals with infinite limits), the function must achieve 
its global maximum at some point strictly within the interior of the interval. 

Thus, the MAP estimate $\hat{s}_j = \arg\max_{s_j} f(s_j)$ 
is guaranteed to satisfy $0 < \hat{s}_j < 1$. The same holds for $g_j$. 
This strict bounding theoretically precludes any point estimate from reaching 
the absorbing states of $0$ or $1$, thereby guaranteeing numerical stability 
during L-BFGS optimization.
\end{proof}

\subsection{Proof of Theorem \ref{thm:strict_concavity}} \label{app:proof_strict_concavity}
\begin{proof}
We compute the first and second derivatives of $\ell(e_j)$ with respect to $e_j$:
\begin{align}
    \frac{\partial \ell(e_j)}{\partial e_j} &= \frac{A_j}{e_j} - \frac{B_j}{1 - e_j} \\
    \frac{\partial^2 \ell(e_j)}{\partial e_j^2} &= -\frac{A_j}{e_j^2} - \frac{B_j}{(1 - e_j)^2}
\end{align}
Since $N_{j1}, S_{j1}, N_{j0}, S_{j0} \ge 0$, and by assumption 
$a > 1 \implies (a-1) > 0$ and $b > 1 \implies (b-1) > 0$, 
it follows that $A_j > 0$ and $B_j > 0$. 
Consequently, for any $e_j \in (0, 1)$:
\begin{equation}
    -\frac{A_j}{e_j^2} - \frac{B_j}{(1 - e_j)^2} < 0
\end{equation}
Because the second derivative is strictly negative across the 
entire domain $(0,1)$, the function is strictly concave. 
Thus, any stationary point found where $\frac{\partial \ell}{\partial e_j} = 0$ 
is guaranteed to be the unique global maximum.
\end{proof}




\newpage
\section{References}
\bibliographystyle{apalike} % Or plain, ieeetr, acm, etc.
\bibliography{references} 

\end{document}
